\documentclass[11pt]{article}
\usepackage[margin=18mm,headheight=22pt]{geometry}
\usepackage{amsmath}
\usepackage{fontspec}
\usepackage{unicode-math}
\setmainfont{NewCM10}[Extension=.otf,UprightFont=*-Book,BoldFont=*-Bold,ItalicFont=*-BookItalic,BoldItalicFont=*-BoldItalic]
\setsansfont{NewCMSans10}[Extension=.otf,UprightFont=*-Book,BoldFont=*-Bold,ItalicFont=*-BookOblique,BoldItalicFont=*-BoldOblique]
\setmathfont{NewCMMath-Book.otf}
\usepackage{xcolor}
\definecolor{FAUblue}{HTML}{04316A}
\definecolor{Paperblue}{HTML}{EEF2F7}
\definecolor{EEIgold}{HTML}{CC9C42}
\usepackage{etoolbox,enumitem,booktabs,tabularx,array,needspace,microtype}
\usepackage[most]{tcolorbox}
\usepackage{titlesec,fancyhdr}
\usepackage[hidelinks]{hyperref}
\setlength{\parindent}{0pt}
\setlength{\parskip}{5pt}
\setlist{itemsep=2pt,topsep=3pt,leftmargin=1.6em,beginpenalty=10000,endpenalty=10000}
\raggedbottom
\clubpenalty=10000 \widowpenalty=10000 \displaywidowpenalty=10000 \interlinepenalty=10000
\setcounter{secnumdepth}{0}
\setcounter{tocdepth}{1}
\pagestyle{fancy}\fancyhf{}
\fancyhead[L]{\small\sffamily\color{FAUblue} From Fundamentals to AI Apocalypse · APDR · Chapters 1+2}
\fancyhead[R]{\small\sffamily Boolean Algebra, Identities, Proofs, and Minimization}
\fancyfoot[C]{\small\thepage}
\titleformat{\section}{\Large\bfseries\sffamily\color{FAUblue}}{}{0pt}{}
\titleformat{\subsection}{\large\bfseries\sffamily\color{FAUblue}}{}{0pt}{}
\titlespacing*{\section}{0pt}{16pt}{6pt}
\preto\section{\Needspace{8\baselineskip}}
\preto\subsection{\Needspace{6\baselineskip}}
\newtcolorbox{equationcard}[1]{colback=Paperblue,colframe=Paperblue,boxrule=0pt,arc=0pt,left=10pt,right=10pt,top=6pt,bottom=6pt,title={#1},fonttitle=\small\bfseries\sffamily,coltitle=FAUblue,before skip=8pt,after skip=8pt}
\newtcolorbox{workedbox}[1]{colback=white,colframe=EEIgold,boxrule=0pt,leftrule=3pt,arc=0pt,left=10pt,top=4pt,bottom=4pt,title={#1},fonttitle=\small\bfseries\sffamily,coltitle=FAUblue,attach title to upper={\par},before skip=8pt,after skip=8pt}
\newtcolorbox{recallbox}[1]{colback=white,colframe=FAUblue,boxrule=0pt,leftrule=3pt,arc=0pt,left=10pt,top=4pt,bottom=4pt,title={#1},fonttitle=\small\bfseries\sffamily,coltitle=FAUblue,attach title to upper={\par},before skip=8pt,after skip=8pt}
\usepackage{environ}
\newsavebox\measurebox
% Typeset the section once invisibly to measure it; start a new page if it will not fit in what is left.
\makeatletter
\newif\ifkeep@measuring
\NewEnviron{keepsection}{%
  \ifkeep@measuring \BODY \else
    \edef\keep@sec{\the\value{section}}\edef\keep@sub{\the\value{subsection}}%
    \setbox\measurebox\vbox{\hsize\linewidth\keep@measuringtrue\def\Needspace##1{}\BODY}%
    \setcounter{section}{\keep@sec}\setcounter{subsection}{\keep@sub}%
    \dimen2=\ht\measurebox \advance\dimen2\dp\measurebox
    \dimen@=\textheight
    \ifdim\pagegoal<\maxdimen
      \dimen@=\pagegoal \advance\dimen@-\pagetotal
      \ifdim\dimen2>\dimen@ \clearpage \dimen@=\textheight \fi
    \fi
    % if it fits, nothing inside may force an early break
    \ifdim\dimen2>\dimen@ \BODY \else \begingroup\def\Needspace##1{}\BODY\endgroup \fi
  \fi}
\makeatother
\newtcolorbox{notebox}{colback=Paperblue,colframe=Paperblue,boxrule=0pt,arc=0pt,left=8pt,right=8pt,top=4pt,bottom=4pt,before skip=6pt,after skip=6pt}
\title{\sffamily\bfseries\color{FAUblue}Boolean Algebra, Identities, Proofs, and Minimization}
\author{Algorithms, Programming, and Data Representation \ensuremath{\cdot} Module 92313 \ensuremath{\cdot} Chapters 1+2\\[2pt]\small From Fundamentals to AI Apocalypse · personal study guide}
\date{}
\begin{document}
\maketitle
\begingroup\footnotesize\setlength{\parskip}{0pt}\tableofcontents\endgroup
\clearpage

\Needspace{8\baselineskip}
{\large This guide develops Boolean logic through algebraic structures, proof tools, function minimization, and decision diagrams.}


\begin{keepsection}
\Needspace{14\baselineskip}
\section{Start from memory}

\Needspace{11\baselineskip}
Answer before reading. Mark each response as sure, unsure, or a guess; uncertainty shows what to revisit.

\Needspace{8\baselineskip}
\begin{enumerate}
\item Name the three operations and two constants in a Boolean algebra.
\item What is the dual of \textit{x} + 0 = \textit{x}?
\item State one absorption law.
\item What two cases does Shannon expansion split on?
\item How many literals remain in a K-map group of 8 cells for a 5-variable function?
\end{enumerate}

\nopagebreak[4]{\small\itshape Answers: see “Start from memory” under Answers at the end.}
\end{keepsection}

\begin{keepsection}
\Needspace{18\baselineskip}
\section{How to use this guide}

\Needspace{15\baselineskip}
Study one short section, close it, then reconstruct its laws or argument from memory. Mix proof, algebra, and diagram questions so you practise choosing the right tool.

\Needspace{12\baselineskip}
\begin{enumerate}
\item \textbf{Day 0:} Try the cold start, then study §§1.1–1.2 and derive duality and De Morgan’s laws.
\item \textbf{Day 1:} Rebuild the NAND proof and normal forms in §1.3; solve mixed practice 1–8.
\item \textbf{Day 3:} Draw K-map groupings from memory and solve practice 9–14.
\item \textbf{About day 7:} Reconstruct Shannon expansion, cofactors, and OBDD reduction; solve practice 15–20.
\item \textbf{About day 14:} Complete both exam-style problems without notes and explain the source checks.
\item \textbf{Before the exam:} Tick only skills you can show cold; return to any section referenced by a missed answer.
\end{enumerate}

\begin{notebox}\textbf{When stuck:} name the definition, write one transformation at a time, and check the result on a small truth table. A guess counts as uncertain until you verify it.\end{notebox}
\end{keepsection}

\begin{keepsection}
\Needspace{17\baselineskip}
\section{How the ideas fit together}

\Needspace{14\baselineskip}
The same Boolean functions can be viewed as algebra, sets, rings, normal forms, maps, or decision diagrams. Each view makes a different task easier.

\Needspace{8\baselineskip}
\textbf{Define} \ensuremath{\cdot} identify the structure

\Needspace{6\baselineskip}
\begin{itemize}
\item 
\Needspace{4\baselineskip}
\textbf{Boolean algebra} \hfill{\small\color{FAUblue}§1.1}

\Needspace{4\baselineskip}
join \ensuremath{\cdot} meet \ensuremath{\cdot} complement

Axioms support equivalent algebraic and lattice views.

\item 
\Needspace{4\baselineskip}
\textbf{Identities and proofs} \hfill{\small\color{FAUblue}§1.2}

\Needspace{4\baselineskip}
duality \ensuremath{\cdot} De Morgan \ensuremath{\cdot} absorption

Derived laws shorten proofs and expressions.

\end{itemize}

\Needspace{7\baselineskip}
\textbf{Represent} \ensuremath{\cdot} choose a form

\Needspace{5\baselineskip}
\begin{itemize}
\item 
\Needspace{4\baselineskip}
\textbf{Complete connectives} \hfill{\small\color{FAUblue}§1.3}

\Needspace{4\baselineskip}
NAND or NOR alone

A complete set can express every Boolean function.

\item 
\Needspace{4\baselineskip}
\textbf{Karnaugh maps} \hfill{\small\color{FAUblue}§1.4}

\Needspace{4\baselineskip}
adjacent cells \ensuremath{\to} fewer literals

Gray-code adjacency exposes variables that can be eliminated.

\end{itemize}

\Needspace{8\baselineskip}
\textbf{Analyse} \ensuremath{\cdot} split and verify

\Needspace{6\baselineskip}
\begin{itemize}
\item 
\Needspace{4\baselineskip}
\textbf{Shannon expansion} \hfill{\small\color{FAUblue}§1.5}

\Needspace{4\baselineskip}
f = x f\textsubscript{1} + x' f\textsubscript{0}

Cofactors lead to decision diagrams and dependency tests.

\item 
\Needspace{4\baselineskip}
\textbf{Verification and hazards} \hfill{\small\color{FAUblue}§1.6}

\Needspace{4\baselineskip}
same function \ensuremath{\cdot} safe transitions

Logical equivalence and physical timing answer different questions.

\end{itemize}

\Needspace{11\baselineskip}
\Needspace{10\baselineskip}
\begin{center}\small\begin{tabularx}{\linewidth}{>{\raggedright\arraybackslash}X>{\raggedright\arraybackslash}X}\toprule
\textbf{Question} & \textbf{Useful view} \\ \midrule
Can two formulas be shown equal? & Use axioms, derived identities, or a truth table. \\
Can a gate set express every function? & Show it can express NOT, AND, and OR. \\
Can a function be represented more compactly? & Try a normal form, a K-map, or a reduced OBDD. \\
\bottomrule\end{tabularx}\end{center}
\end{keepsection}

\begin{keepsection}
\Needspace{6\baselineskip}
\section{1.1 Foundations: one structure, several views}

A Boolean algebra packages truth operations into a system of equations. The same structure can also be understood through order, sets, or arithmetic modulo two.

\begin{keepsection}
\Needspace{14\baselineskip}
\subsection{Operations and axioms}

\Needspace{11\baselineskip}
Write the structure as (B, +, \ensuremath{\cdot}, ', 0, 1). Join (+) and meet (\ensuremath{\cdot}) are binary operations; complement (') is unary. The Huntington-style axioms used in the lecture state commutativity, associativity, distributivity, identities, and complements.

\begin{equationcard}{Core axioms}
\[ \begin{gathered} \mathrm{x}\, +\, \mathrm{y}\, =\, \mathrm{y}\, +\, \mathrm{x} \\ (\mathrm{x}\, +\, \mathrm{y})\, +\, \mathrm{z}\, =\, \mathrm{x}\, +\, (\mathrm{y}\, +\, \mathrm{z}) \\ \mathrm{x}\cdot (\mathrm{y}\, +\, \mathrm{z})\, =\, \mathrm{x}\cdot \mathrm{y}\, +\, \mathrm{x}\cdot \mathrm{z} \\ \mathrm{x}\, +\, 0\, =\, \mathrm{x} \\ \mathrm{x}\cdot 1\, =\, \mathrm{x} \\ \mathrm{x}\, +\, \mathrm{x}'\, =\, 1 \\ \mathrm{x}\cdot \mathrm{x}'\, =\, 0 \end{gathered} \]
{\small The lecture’s axiom list also states the dual commutative and associative laws for meet. Other familiar laws can be derived.}
\end{equationcard}

\Needspace{13\baselineskip}
\textbf{Duality} swaps + with \ensuremath{\cdot} and 0 with 1. Every equation derivable from the axioms has a derivable dual; for example, x + x = x has dual x\ensuremath{\cdot}x = x. A \textbf{distributive complemented bounded lattice} is the same structure in order language: join and meet have bottom 0 and top 1, and each element has a complement.

\begin{workedbox}{Worked proof \ensuremath{\cdot} complements are unique}
Suppose y behaves like a complement of x: x+y=1 and x\ensuremath{\cdot}y=0. Show y=x'.

\Needspace{5\baselineskip}
\begin{enumerate}
\item Use the identity and complement of x: y=y\ensuremath{\cdot}1=y\ensuremath{\cdot}(x+x').
\item Distribute: y=(y\ensuremath{\cdot}x)+(y\ensuremath{\cdot}x')=0+(y\ensuremath{\cdot}x')=y\ensuremath{\cdot}x'.
\item Similarly x'=x'\ensuremath{\cdot}1=x'\ensuremath{\cdot}(x+y)=x'\ensuremath{\cdot}y. Therefore y=y\ensuremath{\cdot}x'=x'\ensuremath{\cdot}y=x'. \hfill{\small\color{FAUblue}§1.1}
\end{enumerate}

The proof relies on the stated complement conditions; it does not assume uniqueness in advance.
\end{workedbox}
\end{keepsection}

\begin{keepsection}
\Needspace{13\baselineskip}
\subsection{Sets, rings, and maps}

\Needspace{10\baselineskip}
The power set P(U), with union as +, intersection as \ensuremath{\cdot}, set complement as ', empty set as 0, and U as 1, is a Boolean algebra. There is also a ring view: XOR is addition, AND is multiplication, and every element is idempotent under multiplication.

\begin{equationcard}{Boolean ring translation}
\[ \begin{gathered} \mathrm{x}\, \text{XOR}\, \mathrm{y}\, =\, (\mathrm{x}\cdot \mathrm{y}')\, +\, (\mathrm{x}'\cdot \mathrm{y}) \\ \mathrm{x}\, ⊙\, \mathrm{y}\, =\, \mathrm{x}\cdot \mathrm{y} \\ \mathrm{x}\, \land \, \mathrm{y}\, =\, \text{xy} \\ \mathrm{x}\, \lor \, \mathrm{y}\, =\, \mathrm{x}\, +\, \mathrm{y}\, +\, \text{xy} \\ \neg \mathrm{x}\, =\, 1\, +\, \mathrm{x} \\ \mathrm{x}²\, =\, \mathrm{x} \\ \mathrm{x}\, +\, \mathrm{x}\, =\, 0 \end{gathered} \]
{\small The power-set example uses symmetric difference for ring addition and intersection for multiplication.}
\end{equationcard}

\Needspace{7\baselineskip}
A homomorphism preserves the operations and constants. Its image is a subalgebra; when it is onto, the lecture’s isomorphism theorem identifies the quotient by its kernel with the target algebra.

\Needspace{11\baselineskip}
\begin{notebox}\textbf{Trap:} duality is a rule for transforming a proved equation; it does not mean complement every variable or reverse every term.\end{notebox}

\begin{recallbox}{Recall check \ensuremath{\cdot} answer without looking}
\Needspace{4\baselineskip}
\begin{enumerate}
\item Which operations and constants appear in the Boolean algebra structure?
\item How does the ring view interpret XOR and AND?
\end{enumerate}

\nopagebreak[4]{\small\itshape Answers: see “1.1 Foundations: one structure, several views” under Answers at the end.}
\end{recallbox}
\end{keepsection}
\end{keepsection}

\begin{keepsection}
\Needspace{15\baselineskip}
\section{1.2 Core identities and proof tools}

\Needspace{12\baselineskip}
Axioms are the starting point; identities are reusable consequences. Proofs become shorter when you name the law that justifies each step.

\Needspace{10\baselineskip}
\begin{center}\small\begin{tabularx}{\linewidth}{>{\raggedright\arraybackslash}X>{\raggedright\arraybackslash}X}\toprule
\textbf{Law} & \textbf{Identity} \\ \midrule
De Morgan & (x+y)'=x'\ensuremath{\cdot}y'(x\ensuremath{\cdot}y)'=x'+y' \\
Absorption & x+x\ensuremath{\cdot}y=xx\ensuremath{\cdot}(x+y)=x \\
Consensus & x\ensuremath{\cdot}y+x'\ensuremath{\cdot}z+y\ensuremath{\cdot}z=x\ensuremath{\cdot}y+x'\ensuremath{\cdot}z \\
Shannon expansion & f=x\ensuremath{\cdot}f\textsubscript{x=1}+x'\ensuremath{\cdot}f\textsubscript{x=0} \\
\bottomrule\end{tabularx}\end{center}

\Needspace{12\baselineskip}
De Morgan’s law follows because x'\ensuremath{\cdot}y' is the complement of x+y: their join is 1 and their meet is 0. The uniqueness-of-complements result from §1.1 then identifies it as (x+y)'. The second law follows by duality.

\begin{workedbox}{Worked simplification \ensuremath{\cdot} remove a redundant condition}
Simplify x+x\ensuremath{\cdot}y using only laws from the lecture.

\Needspace{5\baselineskip}
\begin{enumerate}
\item Factor x with distributivity: x\ensuremath{\cdot}1+x\ensuremath{\cdot}y=x\ensuremath{\cdot}(1+y).
\item Since 1+y=1, this becomes x\ensuremath{\cdot}1.
\item Use the identity law: x\ensuremath{\cdot}1=x. Therefore x+x\ensuremath{\cdot}y=x. \hfill{\small\color{FAUblue}§1.2}
\end{enumerate}

Read the result as: once x is already included, the stricter case x and y adds no new true assignments.
\end{workedbox}

\begin{keepsection}
\Needspace{14\baselineskip}
\subsection{Splitting a function into cases}

\Needspace{11\baselineskip}
The cofactor f\textsubscript{x=1} fixes x to 1; f\textsubscript{x=0} fixes it to 0. Shannon expansion inserts x+x'=1 and distributes f across the two cases. The Boolean difference is the XOR of the cofactors, so it is 1 exactly where changing x changes the output.

\begin{equationcard}{Difference and quantification}
\[ \begin{gathered} \mathrm{D}_\mathrm{x}\, \mathrm{f}\, =\, \mathrm{f}_{\mathrm{x}=1}\, \text{XOR}\, \mathrm{f}_{\mathrm{x}=0} \\ \mathrm{D}_\mathrm{x}(\text{fg})\, =\, \mathrm{f}\cdot \mathrm{D}_\text{xg}\, \text{XOR}\, \mathrm{g}\cdot \mathrm{D}_\text{xf}\, \text{XOR}\, \mathrm{D}_\text{xf}\cdot \mathrm{D}_\text{xg} \\ (\exists \mathrm{x})\mathrm{f}\, =\, \mathrm{f}_{\mathrm{x}=1}\, +\, \mathrm{f}_{\mathrm{x}=0} \\ (\forall \mathrm{x})\mathrm{f}\, =\, \mathrm{f}_{\mathrm{x}=1}\cdot \mathrm{f}_{\mathrm{x}=0} \end{gathered} \]
{\small Existential quantification asks whether either case works; universal quantification asks whether both cases work. Both remove dependence on x.}
\end{equationcard}

\Needspace{11\baselineskip}
\begin{notebox}\textbf{Memory hook:} for existential, OR the two cofactors; for universal, AND them. Check the question’s quantifier before simplifying.\end{notebox}

\begin{recallbox}{Recall check \ensuremath{\cdot} answer without looking}
\Needspace{4\baselineskip}
\begin{enumerate}
\item Why can the complement uniqueness theorem prove De Morgan’s first law?
\item What does a Boolean difference of 0 tell you about x?
\end{enumerate}

\nopagebreak[4]{\small\itshape Answers: see “1.2 Core identities and proof tools” under Answers at the end.}
\end{recallbox}
\end{keepsection}
\end{keepsection}

\begin{keepsection}
\Needspace{6\baselineskip}
\section{1.3 Completeness, normal forms, and minimization}

A representation tells you how to build a function. Completeness says whether a chosen set of connectives can represent every function; minimization asks how small a representation can be.

\begin{keepsection}
\Needspace{13\baselineskip}
\subsection{Build every function from a small gate set}

\Needspace{10\baselineskip}
A connective set is functionally complete if it can express every Boolean function. NAND alone is complete because it can express NOT, AND, and OR. NOR alone is complete by the dual argument.

\begin{equationcard}{NAND construction}
\[ \begin{gathered} \mathrm{x}\, ↑\, \mathrm{y}\, =\, (\mathrm{x}\cdot \mathrm{y})' \\ \neg \mathrm{x}\, =\, \mathrm{x}\, ↑\, \mathrm{x} \\ \mathrm{x}\cdot \mathrm{y}\, =\, (\mathrm{x}\, ↑\, \mathrm{y})\, ↑\, (\mathrm{x}\, ↑\, \mathrm{y}) \\ \mathrm{x}+\mathrm{y}\, =\, (\mathrm{x}\, ↑\, \mathrm{x})\, ↑\, (\mathrm{y}\, ↑\, \mathrm{y}) \end{gathered} \]
{\small Once NOT, AND, and OR are available, any formula built from those connectives can be expressed with NAND.}
\end{equationcard}

Post’s lattice classifies families of Boolean functions closed under composition. The lecture also contrasts monotone functions, which preserve order, with affine functions, which are parity expressions over F\textsubscript{2}; a function that is both monotone and affine is constant.
\end{keepsection}

\begin{keepsection}
\Needspace{15\baselineskip}
\subsection{Write and minimize formulas}

\Needspace{12\baselineskip}
Every Boolean function has a disjunctive normal form (OR of minterms for true rows) and a conjunctive normal form (AND of maxterms that exclude false rows). These forms are systematic, though a minimum exact two-level DNF or CNF search is NP-complete, so practical methods use structure and heuristics.

\begin{workedbox}{Worked example \ensuremath{\cdot} XOR as DNF and CNF}
Let f(A,B)=1 on rows 01 and 10, and 0 on rows 00 and 11.

\Needspace{5\baselineskip}
\begin{enumerate}
\item For the true rows, the minterms are A'B and AB', so the DNF is A'B+AB'.
\item The false rows 00 and 11 give maxterms (A+B) and (A'+B'), so the CNF is (A+B)(A'+B').
\item Both forms describe XOR: the output is 1 exactly when the inputs differ. \hfill{\small\color{FAUblue}§1.3}
\end{enumerate}
\end{workedbox}

\Needspace{6\baselineskip}
Quine–McCluskey combines minterms that differ in one literal, repeats until prime implicants remain, then selects a cover. Its steps preserve coverage: every combined implicant covers only true assignments, and every listed true minterm remains covered.

\Needspace{10\baselineskip}
\begin{notebox}\textbf{Trap:} a complete gate set guarantees expressiveness, not a small or fast circuit.\end{notebox}

\begin{recallbox}{Recall check \ensuremath{\cdot} answer without looking}
\Needspace{4\baselineskip}
\begin{enumerate}
\item How does NAND express NOT, AND, and OR?
\item What does DNF list, and why is finding a minimum exact form difficult?
\end{enumerate}

\nopagebreak[4]{\small\itshape Answers: see “1.3 Completeness, normal forms, and minimization” under Answers at the end.}
\end{recallbox}
\end{keepsection}
\end{keepsection}

\begin{keepsection}
\Needspace{11\baselineskip}
\section{1.4 Karnaugh maps remove changing variables}

\Needspace{8\baselineskip}
A K-map lays truth-table rows out so neighboring cells differ in one variable. Grouping adjacent 1s makes that changing variable disappear from the corresponding implicant.

\Needspace{11\baselineskip}
For n variables, a map has 2\textsuperscript{n} cells in Gray-code order. A legal group contains 2\textsuperscript{k} adjacent cells, with wraparound adjacency allowed. Each grouped variable that changes across the group is eliminated; the implicant keeps n\ensuremath{-}k literals. Use don’t-care rows as either 0 or 1 when that makes a useful group.

\begin{equationcard}{K-map group rule}
\[ \begin{gathered} \text{Map}\, \text{size}:\, 2^{\mathrm{n}}\, \text{cells} \\ \text{Group}\, \text{size}:\, 2^{\mathrm{k}}\, \text{cells} \\ \text{Remaining}\, \text{literals}:\, \mathrm{n}-\mathrm{k} \end{gathered} \]
{\small Cover every required 1 at least once. Do not include a required 0. A don’t-care can be included or ignored.}
\end{equationcard}

\begin{workedbox}{Worked example \ensuremath{\cdot} simplify using an explicit bit order}
Use variable order A,B,C with A as the most significant bit. Suppose f=1 for minterms 2, 3, 6, and 7.

\Needspace{5\baselineskip}
\begin{enumerate}
\item These rows are 010, 011, 110, and 111, so B=1 in every true row.
\item On the map, they form a group of four. A and C vary; B stays fixed.
\item The group leaves one literal: f=B. Verify that B=0 rows map to 0 and B=1 rows map to 1. \hfill{\small\color{FAUblue}§1.4}
\end{enumerate}
\end{workedbox}

\Needspace{10\baselineskip}
\begin{notebox}\textbf{Trap:} write the variable bit order beside minterm indices. A minterm number has no meaning without a variable order.\end{notebox}

\begin{recallbox}{Recall check \ensuremath{\cdot} answer without looking}
\Needspace{4\baselineskip}
\begin{enumerate}
\item How many literals remain after grouping 8 cells on a 5-variable map?
\item What is the safe way to use a don’t-care cell?
\end{enumerate}

\nopagebreak[4]{\small\itshape Answers: see “1.4 Karnaugh maps remove changing variables” under Answers at the end.}
\end{recallbox}
\end{keepsection}

\begin{keepsection}
\Needspace{12\baselineskip}
\section{1.5 Shannon views, OBDDs, and algebraic normal form}

\Needspace{9\baselineskip}
Shannon expansion recursively splits a function into simpler cases. Reusing the same split structure leads to decision diagrams; a different basis expresses the function as parity over F\textsubscript{2}.

\begin{equationcard}{Shannon expansion}
\[ \begin{gathered} \mathrm{f}\, =\, \mathrm{x}\cdot \mathrm{f}_{\mathrm{x}=1}\, +\, \mathrm{x}'\cdot \mathrm{f}_{\mathrm{x}=0} \end{gathered} \]
{\small The selected input x chooses between the two cofactors, like a 2-to-1 multiplexer.}
\end{equationcard}

\begin{workedbox}{Worked example \ensuremath{\cdot} expand on A}
For f(A,B)=A\ensuremath{\cdot}B+A', compute the two cofactors and reconstruct f.

\Needspace{5\baselineskip}
\begin{enumerate}
\item Set A=1: f\textsubscript{A=1}=B+0=B. Set A=0: f\textsubscript{A=0}=0+1=1.
\item Substitute: f=A\ensuremath{\cdot}B+A'\ensuremath{\cdot}1.
\item The expansion matches the starting expression and separates the A=1 and A=0 cases. \hfill{\small\color{FAUblue}§1.5}
\end{enumerate}
\end{workedbox}

\Needspace{10\baselineskip}
An ordered binary decision diagram (OBDD) applies Shannon expansion in a fixed variable order. Reduction merges isomorphic subgraphs and removes tests whose two branches are identical. For one fixed order, the reduced OBDD is unique; a different variable order can produce a different diagram size.

\Needspace{11\baselineskip}
The algebraic normal form (ANF) writes f as XORs of monomials over F\textsubscript{2}. Every Boolean function has a unique ANF; its coefficients can be computed by a Möbius transform. Uniqueness follows because the monomials form a basis for all truth functions on n variables. Unlike DNF, ANF uses XOR as its addition operation.

\begin{equationcard}{ANF shape}
\[ \begin{gathered} \mathrm{f}\, \text{is}\, \text{the}\, \text{XOR}\, \text{of}\, \mathrm{a}_{\mathrm{S}}\cdot ∏_{\mathrm{i}\in \mathrm{S}}\mathrm{x}_{\mathrm{i}}\, \text{over}\, \mathrm{S}⊆[\mathrm{n}] \end{gathered} \]
{\small Each coefficient a\textsubscript{S} is 0 or 1. The empty product is 1; the XOR collects the selected monomials.}
\end{equationcard}

\Needspace{10\baselineskip}
\begin{notebox}\textbf{Trap:} an OBDD is only canonical after you fix the variable order.\end{notebox}

\begin{recallbox}{Recall check \ensuremath{\cdot} answer without looking}
\Needspace{4\baselineskip}
\begin{enumerate}
\item How does a Shannon expansion combine its two cofactors?
\item What conditions make an OBDD reduced and canonical?
\end{enumerate}

\nopagebreak[4]{\small\itshape Answers: see “1.5 Shannon views, OBDDs, and algebraic normal form” under Answers at the end.}
\end{recallbox}
\end{keepsection}

\begin{keepsection}
\Needspace{6\baselineskip}
\section{1.6 Verification, hazards, and wrap-up}

Equivalent formulas can behave differently while a real circuit switches. Verification must check both the Boolean function and, where needed, the physical transitions.

\begin{keepsection}
\Needspace{9\baselineskip}
\subsection{Static-1 hazards}

\begin{equationcard}{Consensus relation}
\[ \begin{gathered} \text{xy}\, +\, \mathrm{x}'\mathrm{z}\, +\, \text{yz}\, =\, \text{xy}\, +\, \mathrm{x}'\mathrm{z} \end{gathered} \]
{\small The term yz is logically redundant but can preserve coverage during a circuit transition.}
\end{equationcard}

\Needspace{13\baselineskip}
A static-1 hazard occurs when the output should stay 1 for a small input change but a particular two-level sum-of-products implementation briefly drops to 0 because its paths have different delays. If two product terms hand off across one changing literal, their consensus term supplies overlapping coverage.

\begin{workedbox}{Worked example \ensuremath{\cdot} add the consensus term}
For f=xy+x'z, the consensus term is yz.

\Needspace{6\baselineskip}
\begin{enumerate}
\item The terms xy and x'z differ on x; both can be 1 at the boundary when y=z=1.
\item Add yz. The consensus theorem gives xy+x'z+yz=xy+x'z, so the Boolean function is unchanged.
\item In a two-level sum-of-products implementation, yz keeps the output covered during the handoff and removes the static-1 hazard described in the lecture. \hfill{\small\color{FAUblue}§1.6}
\end{enumerate}
\end{workedbox}
\end{keepsection}

\begin{keepsection}
\Needspace{12\baselineskip}
\subsection{Soundness and completeness}

\Needspace{9\baselineskip}
Equational reasoning is sound when every derived equation is valid in every Boolean algebra model. It is complete for tautological equivalence when every such equivalence can be derived using the calculus. The lecture’s wrap-up ties this proof view to the alternate ring, set, map, and decision-diagram representations. Its further-reading slide points to Stone on representation, Post on completeness, and Bryant on decision diagrams.

\Needspace{11\baselineskip}
\begin{notebox}\textbf{Sign check:} Boolean equivalence says outputs agree for stable inputs; it does not guarantee zero propagation delay in hardware.\end{notebox}

\begin{recallbox}{Recall check \ensuremath{\cdot} answer without looking}
\Needspace{4\baselineskip}
\begin{enumerate}
\item Why can a redundant consensus term help a physical circuit?
\item What is the difference between soundness and completeness?
\end{enumerate}

\nopagebreak[4]{\small\itshape Answers: see “1.6 Verification, hazards, and wrap-up” under Answers at the end.}
\end{recallbox}
\end{keepsection}
\end{keepsection}

\begin{keepsection}
\Needspace{19\baselineskip}
\section{Formula map}

\Needspace{10\baselineskip}
\begin{center}\small\begin{tabularx}{\linewidth}{>{\raggedright\arraybackslash}X>{\raggedright\arraybackslash}X>{\raggedright\arraybackslash}X}\toprule
\textbf{Idea} & \textbf{Relation} & \textbf{Use} \\ \midrule
Duality & Swap + with \ensuremath{\cdot} and 0 with 1 & Derive the matching law. \\
De Morgan & (x+y)'=x'\ensuremath{\cdot}y'(x\ensuremath{\cdot}y)'=x'+y' & Move a complement across an operation. \\
Absorption & x+x\ensuremath{\cdot}y=x & Remove a term already covered by x. \\
Consensus & xy+x'z+yz=xy+x'z & Remove a logically redundant term; add it back to cover a hazard. \\
Shannon & f=xf\textsubscript{1}+x'f\textsubscript{0} & Split on a variable. \\
Boolean difference & D\_x f=f\textsubscript{1} XOR f\textsubscript{0} & Test whether x can change the result. \\
K-map group & 2\textsuperscript{k} cells leave n\ensuremath{-}k literals & Eliminate variables that change. \\
NAND & x↑y=(xy)' & Build NOT, AND, and OR from one gate type. \\
\bottomrule\end{tabularx}\end{center}
\end{keepsection}

\begin{keepsection}
\Needspace{20\baselineskip}
\section{Derive it cold}

\Needspace{17\baselineskip}
Cover the right column and rebuild each chain using only earlier definitions and laws.

\Needspace{15\baselineskip}
\begin{enumerate}
\item \textbf{Duality:} x+x=x \ensuremath{\to} x\ensuremath{\cdot}x=x. \hfill{\small\color{FAUblue}§1.1}
\item \textbf{De Morgan:} show x'y' complements x+y, then use uniqueness. \hfill{\small\color{FAUblue}§1.2}
\item \textbf{Absorption:} x+xy=x(1+y)=x. \hfill{\small\color{FAUblue}§1.2}
\item \textbf{Shannon:} f=f(x+x')=xf+x'f, then read each term as a cofactor. \hfill{\small\color{FAUblue}§1.2}
\item \textbf{NAND completeness:} derive NOT, AND, and OR, then cite functional completeness of those connectives. \hfill{\small\color{FAUblue}§1.3}
\item \textbf{K-map grouping:} pair adjacent cells and show the changing variable disappears. \hfill{\small\color{FAUblue}§1.4}
\item \textbf{OBDD uniqueness:} explain why the fixed order plus reduction rules determine the same graph. \hfill{\small\color{FAUblue}§1.5}
\item \textbf{Hazard cover:} expand yz with x+x' and show it is absorbed by xy+x'z. \hfill{\small\color{FAUblue}§1.6}
\end{enumerate}
\end{keepsection}

\begin{keepsection}
\Needspace{16\baselineskip}
\section{Explain it}

\Needspace{13\baselineskip}
Answer each “why” or “when” in your own words before opening the model response.

\Needspace{11\baselineskip}
\begin{enumerate}
\item Why does duality save proof work?
\item Why is uniqueness of complements useful for proving De Morgan’s law?
\item When is an apparently shorter Boolean expression unsafe to use as a circuit simplification?
\item Why can NAND alone express every Boolean function?
\item When should a don’t-care be included in a K-map group?
\item Why does an OBDD’s variable order matter even when its reduced form is canonical?
\item Why does a K-map group eliminate the variables that change?
\end{enumerate}

\nopagebreak[4]{\small\itshape Answers: see “Explain it” under Answers at the end.}
\end{keepsection}

\begin{keepsection}
\Needspace{17\baselineskip}
\section{Source checks}

\Needspace{14\baselineskip}
These two items were checked against the rendered slides. The guide states its indexing convention where it affects the result.

\Needspace{11\baselineskip}
\begin{itemize}
\item \textbf{Slide 35 \ensuremath{\cdot} minterm order not stated:} the slide gives \ensuremath{\Sigma}m(1,3,5,7) and concludes f=B, but it does not label the variable bit order. Under the common A,B,C order with A most significant, that set describes C; the slide’s conclusion matches a different bit order with B least significant. This guide states the convention explicitly in its worked example.
\item \textbf{Slide 36 \ensuremath{\cdot} expression conflicts with listed rows:} with W,X,Y,Z in that bit order, the listed 1s \{1,3,7,11,15\} and don’t-cares \{5,9,13\} are all odd indices, while every even index is 0; the list therefore simplifies to f=Z. The rendered slide’s printed expression includes W'Z', which is 1 on some even rows. This guide follows the listed rows and flags the mismatch for checking against the lecturer’s source.
\end{itemize}
\end{keepsection}

\clearpage

\begin{keepsection}
\Needspace{10\baselineskip}
\section{Mixed practice}

\Needspace{7\baselineskip}
Try all 20 questions in order without notes. The topics are interleaved; write your confidence beside each answer before checking the hidden key.

\begin{enumerate}
\item State the dual of x+0=x. \hfill{\small\color{FAUblue}§1.1}
\item Simplify a+a\ensuremath{\cdot}b. Name the law. \hfill{\small\color{FAUblue}§1.2}
\item Using NAND ↑, write NOT x. \hfill{\small\color{FAUblue}§1.3}
\item How many cells are on a 5-variable K-map? \hfill{\small\color{FAUblue}§1.4}
\item For f(x,y)=xy+x', find f\textsubscript{x=1} and f\textsubscript{x=0}. \hfill{\small\color{FAUblue}§1.5}
\item What consensus term can be added to xy+x'z to cover a static-1 handoff? \hfill{\small\color{FAUblue}§1.6}
\item In a Boolean ring, calculate 1+1 and 1\ensuremath{\cdot}1. \hfill{\small\color{FAUblue}§1.1}
\item Rewrite (a+b)' using De Morgan’s law. \hfill{\small\color{FAUblue}§1.2}
\item For f(A,B)=1 on rows 01 and 10 only, give one DNF. \hfill{\small\color{FAUblue}§1.3}
\item A group of 8 cells is used on a 5-variable K-map. How many literals remain? \hfill{\small\color{FAUblue}§1.4}
\item Expand f(A,B)=AB+A' on A by listing the two cofactors. \hfill{\small\color{FAUblue}§1.5}
\item Distinguish soundness from completeness in equational reasoning. \hfill{\small\color{FAUblue}§1.6}
\item Give the dual of x+x'=1. \hfill{\small\color{FAUblue}§1.1}
\item What is Dx/Dx? \hfill{\small\color{FAUblue}§1.2}
\item A function is both monotone and affine. What does the lecture say it must be? \hfill{\small\color{FAUblue}§1.3}
\item With WXYZ as the bit order, use the slide 36 listed rows: 1s at 1,3,7,11,15 and don’t-cares at 5,9,13. What is the simplest function? \hfill{\small\color{FAUblue}§1.4}
\item Write A OR B in algebraic normal form over F\textsubscript{2}. \hfill{\small\color{FAUblue}§1.5}
\item Does adding yz change the Boolean function xy+x'z? \hfill{\small\color{FAUblue}§1.6}
\item Why can exact minimum two-level DNF search require a heuristic in practice? \hfill{\small\color{FAUblue}§1.3}
\item What does a Boolean difference of 0 with respect to x imply? \hfill{\small\color{FAUblue}§1.2}
\end{enumerate}

\nopagebreak[4]{\small\itshape Answers: see “Mixed practice” under Answers at the end.}
\end{keepsection}

\begin{keepsection}
\Needspace{6\baselineskip}
\section{Exam-style problems}

Allow about 20 minutes per problem. Show each identity used and check your final expression against the truth conditions.

\begin{keepsection}
\Needspace{12\baselineskip}
\subsection{1 \ensuremath{\cdot} Prove, represent, and implement}

\Needspace{9\baselineskip}
For f(A,B,C)=A'B+AB'C, use the Boolean algebra and representation ideas from §§1.1–1.3.

\Needspace{7\baselineskip}
\begin{enumerate}
\item Use Shannon expansion on A to write the two cofactors and reconstruct f.
\item Write a NAND-only implementation for each product term and the outer OR, using NAND as the only gate.
\item State whether the given expression is a DNF and explain how you would verify equivalence to your NAND version.
\end{enumerate}

\nopagebreak[4]{\small\itshape Answers: see “Exam-style problems” under Answers at the end.}
\end{keepsection}

\begin{keepsection}
\Needspace{11\baselineskip}
\subsection{2 \ensuremath{\cdot} Minimize, then check a circuit transition}

\Needspace{8\baselineskip}
Let f(A,B,C)=1 for minterms 2,3,6,7 and 0 elsewhere. Use A,B,C as the bit order with A most significant.

\Needspace{5\baselineskip}
\begin{enumerate}
\item Group the true cells on a 3-variable K-map and derive the minimal expression.
\item Expand your result on A using Shannon’s theorem.
\item Explain what a static-1 hazard is and when the consensus theorem helps with it.
\end{enumerate}

\nopagebreak[4]{\small\itshape Answers: see “Exam-style problems” under Answers at the end.}
\end{keepsection}
\end{keepsection}

\begin{keepsection}
\Needspace{10\baselineskip}
\section{Self-check}

\Needspace{7\baselineskip}
Use this list in a later session, at least a day after learning the material. Tick a skill only when you can demonstrate it without notes.

\begin{itemize}
\item $\mdlgwhtsquare$\ State the Boolean operations, constants, and core axioms. (1, 7)
\item $\mdlgwhtsquare$\ Apply duality to a proved equation. (1, 13)
\item $\mdlgwhtsquare$\ Prove complement uniqueness and use it for De Morgan. (Derive it cold)
\item $\mdlgwhtsquare$\ Simplify with absorption and consensus. (2, 18)
\item $\mdlgwhtsquare$\ Compute cofactors, Shannon expansion, and Boolean difference. (5, 11, 14, 20)
\item $\mdlgwhtsquare$\ Distinguish existential and universal quantification.
\item $\mdlgwhtsquare$\ Translate XOR and AND in the Boolean ring view. (7, 17)
\item $\mdlgwhtsquare$\ Build NOT, AND, and OR from NAND. (3)
\item $\mdlgwhtsquare$\ Create DNF and CNF from truth conditions. (9)
\item $\mdlgwhtsquare$\ Explain why minimum exact forms can be computationally hard. (19)
\item $\mdlgwhtsquare$\ Group K-map cells and count the remaining literals. (4, 10, 16)
\item $\mdlgwhtsquare$\ State the variable order used by minterm indices. (16)
\item $\mdlgwhtsquare$\ Explain OBDD reduction and fixed-order canonicity.
\item $\mdlgwhtsquare$\ Write a simple function in algebraic normal form. (17)
\item $\mdlgwhtsquare$\ Distinguish Boolean equivalence from hazard-free switching. (6, 18)
\item $\mdlgwhtsquare$\ Rebuild the derivations from their definitions.
\item $\mdlgwhtsquare$\ Complete both multi-part problems and justify each step.
\end{itemize}
\end{keepsection}


\clearpage
\section{Answers}
\begin{keepsection}\subsection{Start from memory}

\Needspace{8\baselineskip}
\begin{enumerate}
\item Join (+), meet (\ensuremath{\cdot}), complement ('), and constants 0 and 1. \hfill{\small\color{FAUblue}§1.1}
\item \textit{x} \ensuremath{\cdot} 1 = \textit{x}, swapping + with \ensuremath{\cdot} and 0 with 1. \hfill{\small\color{FAUblue}§1.1}
\item \textit{x} + \textit{x}\ensuremath{\cdot}\textit{y} = \textit{x}, or its dual. \hfill{\small\color{FAUblue}§1.2}
\item It fixes the selected variable first to 1 and then to 0, combining the two cofactors. \hfill{\small\color{FAUblue}§1.5}
\item Two literals remain: a group of 2\textsuperscript{3} eliminates 3 of 5 variables. \hfill{\small\color{FAUblue}§1.4}
\end{enumerate}

\end{keepsection}
\begin{keepsection}\subsection{1.1 Foundations: one structure, several views}

\Needspace{4\baselineskip}
\begin{enumerate}
\item Binary join +, binary meet \ensuremath{\cdot}, unary complement ', and constants 0 and 1. \hfill{\small\color{FAUblue}§1.1}
\item XOR is ring addition and AND is ring multiplication; each element satisfies x²=x. \hfill{\small\color{FAUblue}§1.1}
\end{enumerate}

\end{keepsection}
\begin{keepsection}\subsection{1.2 Core identities and proof tools}

\Needspace{5\baselineskip}
\begin{enumerate}
\item Show x'\ensuremath{\cdot}y' has join 1 and meet 0 with x+y; it is therefore the unique complement of x+y. \hfill{\small\color{FAUblue}§1.2}
\item The two cofactors agree everywhere, so flipping x never changes f; f does not depend on x. \hfill{\small\color{FAUblue}§1.2}
\end{enumerate}

\end{keepsection}
\begin{keepsection}\subsection{1.3 Completeness, normal forms, and minimization}

\Needspace{5\baselineskip}
\begin{enumerate}
\item NOT x=x↑x; AND x\ensuremath{\cdot}y=(x↑y)↑(x↑y); OR x+y=(x↑x)↑(y↑y). \hfill{\small\color{FAUblue}§1.3}
\item DNF ORs minterms for true rows. Exact minimum two-level minimization is NP-complete, so structured methods and heuristics matter. \hfill{\small\color{FAUblue}§1.3}
\end{enumerate}

\end{keepsection}
\begin{keepsection}\subsection{1.4 Karnaugh maps remove changing variables}

\Needspace{4\baselineskip}
\begin{enumerate}
\item Five minus three changing variables leaves two literals. \hfill{\small\color{FAUblue}§1.4}
\item Choose 0 or 1 if useful, but never use a don’t-care to cover a required zero. \hfill{\small\color{FAUblue}§1.4}
\end{enumerate}

\end{keepsection}
\begin{keepsection}\subsection{1.5 Shannon views, OBDDs, and algebraic normal form}

\Needspace{5\baselineskip}
\begin{enumerate}
\item Multiply the x=1 cofactor by x, the x=0 cofactor by x', then OR the two cases. \hfill{\small\color{FAUblue}§1.5}
\item Fix a variable order, merge isomorphic subgraphs, and remove redundant tests; uniqueness is for that fixed order. \hfill{\small\color{FAUblue}§1.5}
\end{enumerate}

\end{keepsection}
\begin{keepsection}\subsection{1.6 Verification, hazards, and wrap-up}

\Needspace{6\baselineskip}
\begin{enumerate}
\item It preserves overlapping coverage while unequal path delays could otherwise create a brief 1-to-0 glitch. \hfill{\small\color{FAUblue}§1.6}
\item Soundness means every proved equation is valid. Completeness means every tautological equivalence is derivable. \hfill{\small\color{FAUblue}§1.6}
\end{enumerate}

\end{keepsection}
\begin{keepsection}\subsection{Explain it}

\Needspace{22\baselineskip}
\begin{enumerate}
\item The duality principle preserves derivability when + and \ensuremath{\cdot}, and 0 and 1, are exchanged. Once one equation is proved, its dual follows from the same axiom system without repeating the argument.
\item It is enough to verify that the proposed expression has join 1 and meet 0 with the original expression. Those two properties make it a complement, and uniqueness identifies it with the negation.
\item Logical equivalence alone does not account for unequal propagation delays. A two-level sum-of-products implementation may glitch during a transition, so a consensus term can be needed even when it is algebraically redundant.
\item NAND can produce NOT by tying its inputs together. It can then produce AND and OR, and those connectives can express every Boolean function.
\item Include it when it helps make a larger legal group and removes literals. Ignore it when it does not help; it is not a required 1 and must never force a group through a required 0.
\item Canonicity holds for a fixed order. Changing the order changes the sequence of Shannon splits and can therefore change the reduced diagram’s shape and size.
\item The group contains every combination of the changing variables while the remaining variables stay fixed. Combining those cases with Boolean identities removes the changing literals, leaving only the fixed ones.
\end{enumerate}

\end{keepsection}
\begin{keepsection}\subsection{Mixed practice}

\begin{enumerate}
\item x\ensuremath{\cdot}1=x, swapping join with meet and 0 with 1. \hfill{\small\color{FAUblue}§1.1}
\item a+a\ensuremath{\cdot}b=a by absorption. \hfill{\small\color{FAUblue}§1.2}
\item x'=x↑x because NAND of x with itself negates x. \hfill{\small\color{FAUblue}§1.3}
\item 2\textsuperscript{5}=32 cells, one per input assignment. \hfill{\small\color{FAUblue}§1.4}
\item f\textsubscript{x=1}=y+0=y; f\textsubscript{x=0}=0+1=1. \hfill{\small\color{FAUblue}§1.5}
\item yz, giving xy+x'z+yz. The added term preserves the function and overlaps the handoff. \hfill{\small\color{FAUblue}§1.6}
\item 1+1=0 by characteristic 2; 1\ensuremath{\cdot}1=1 by the multiplicative identity. \hfill{\small\color{FAUblue}§1.1}
\item (a+b)'=a'b'. \hfill{\small\color{FAUblue}§1.2}
\item A'B+AB', the two minterms for rows 01 and 10. \hfill{\small\color{FAUblue}§1.3}
\item Two literals: n\ensuremath{-}k=5\ensuremath{-}3=2. \hfill{\small\color{FAUblue}§1.4}
\item f\textsubscript{A=1}=B and f\textsubscript{A=0}=1, so f=AB+A'\ensuremath{\cdot}1. \hfill{\small\color{FAUblue}§1.5}
\item Soundness means every derivable equation is valid; completeness means every tautological equivalence is derivable. \hfill{\small\color{FAUblue}§1.6}
\item x\ensuremath{\cdot}x'=0, the dual of x+x'=1. \hfill{\small\color{FAUblue}§1.1}
\item 1: the cofactors of x with respect to itself are 1 and 0, whose XOR is 1. \hfill{\small\color{FAUblue}§1.2}
\item It must be constant, as stated in the lecture. \hfill{\small\color{FAUblue}§1.3}
\item Every odd index is either 1 or a don’t-care and every even index is 0, so f=Z under the stated order. \hfill{\small\color{FAUblue}§1.4}
\item A OR B = A XOR B XOR AB. Check rows 00, 01, 10, and 11 to verify the outputs 0,1,1,1. \hfill{\small\color{FAUblue}§1.5}
\item No. The consensus theorem says xy+x'z+yz=xy+x'z. \hfill{\small\color{FAUblue}§1.6}
\item Exact two-level minimization is NP-complete, so the search can grow too costly; structured methods and heuristics help. \hfill{\small\color{FAUblue}§1.3}
\item The cofactors agree everywhere, so changing x never changes the function. \hfill{\small\color{FAUblue}§1.2}
\end{enumerate}

\end{keepsection}
\begin{keepsection}\subsection{Exam-style problems}

\Needspace{9\baselineskip}
\begin{enumerate}
\item For A=1, f=B'C; for A=0, f=B. Therefore f=A\ensuremath{\cdot}(B'C)+A'\ensuremath{\cdot}B, the original function after reordering its terms. \hfill{\small\color{FAUblue}§1.2} \hfill{\small\color{FAUblue}§1.5}
\item Make complements with x'=x↑x. For p=A'B, compute u=A'↑B=(A'B)', then p=u↑u. For q=AB'C, compute v=A↑B'=(AB')', r=v↑v=AB', s=r↑C=(AB'C)', then q=s↑s. Finally f=(p↑p)↑(q↑q), which is p OR q by De Morgan. \hfill{\small\color{FAUblue}§1.3}
\item Yes: it is an OR of product terms. A truth table or algebraic rewriting can show that the gate expression produces the same value on every assignment. \hfill{\small\color{FAUblue}§1.3}
\end{enumerate}

\end{keepsection}
\begin{keepsection}\subsection{Exam-style problems}

\Needspace{9\baselineskip}
\begin{enumerate}
\item The rows 010, 011, 110, and 111 all have B=1; they form a group of four. A and C vary, so f=B. \hfill{\small\color{FAUblue}§1.4}
\item B is independent of A, so both A cofactors equal B. Shannon gives f=AB+A'B=B(A+A')=B. \hfill{\small\color{FAUblue}§1.2} \hfill{\small\color{FAUblue}§1.5}
\item A static-1 hazard is a brief drop to 0 when the intended output stays 1 during a small input change. In a two-level sum-of-products implementation, an added consensus term can maintain overlapping coverage without changing the Boolean function. \hfill{\small\color{FAUblue}§1.6}
\end{enumerate}

\end{keepsection}

\vfill{\small Study guide based on Bernhard Kainz, “Boolean Algebra, Identities, Proofs, and Minimization: An advanced mathematical perspective for CS undergraduates,” supplied as APDR lecture01 - extended.pdf (48 slides, dated 5 October 2025). Written in my own words; no lecture pages or figures are reproduced. Source inconsistencies checked against rendered slides 35–36 are described above. Not an official FAU publication; check course details against the current materials.}

\end{document}
